\chapter{Glossaire}

\begin{description}
    \item[Cahier des Charges] Document synthétique (3 à 5~pages dans SRS) décrivant les
    besoins fonctionnels et non fonctionnels essentiels d'un projet.
    \item[Knowledge Graph] Modèle interne de SRS structurant les réponses de l'entretien
    par domaine et par dépendances entre concepts.
    \item[RBAC] \emph{Role-Based Access Control} — contrôle d'accès basé sur les rôles.
    \item[STRIDE] Méthode de modélisation des menaces (Spoofing, Tampering,
    Repudiation, Information disclosure, Denial of service, Elevation of privilege).
    \item[DREAD] Méthode de pondération des risques (Damage, Reproducibility,
    Exploitability, Affected users, Discoverability).
    \item[PIA] \emph{Privacy Impact Assessment} — analyse d'impact relative à la
    protection des données.
    \item[MVP] \emph{Minimum Viable Product} — version minimale mais fonctionnelle d'un
    produit.
    \item[ODD] Objectifs de Développement Durable (Nations Unies).
    \item[ESG] Environnement, Social, Gouvernance — cadre d'indicateurs de durabilité.
    \item[SRS] Secure Requirements Studio — l'outil objet de ce rapport.
\end{description}
